\documentclass[12pt,a4paper]{article}

\usepackage{fontspec}
\setmainfont{Latin Modern Roman}

\usepackage[margin=1.8cm]{geometry}
\usepackage{setspace}
\usepackage{parskip}
\usepackage{fancyhdr}

\setstretch{1.08}
\pagestyle{fancy}
\fancyhf{}
\renewcommand{\headrulewidth}{0pt}
\renewcommand{\footrulewidth}{0pt}

\newcommand{\slidenum}[1]{%
  \par\vspace{0.55cm}
  {\bfseries\large #1.}\quad
}

\begin{document}

\slidenum{1}
Buongiorno, sono Samuele Schiavi e presento la mia tesi di laurea triennale in Fisica che affronta il tema della simulazione quantistica.

\slidenum{2}
Il contesto è quello dei magneti molecolari: sistemi basati su spin 1/2 accoppiati, che costituiscono piattaforme privilegiate per lo studio di fenomeni quantistici la cui simulazione classica diventa rapidamente intrattabile al crescere del numero di spin coinvolti. Lo spazio di Hilbert cresce come 2 alla N e attorno alla cinquantina la diagonalizzazione esatta esaurisce le risorse dei supercomputer attuali. I computer quantistici offrono un percorso naturale per aggirare questo limite, perché ogni spin 1/2 si mappa direttamente su un qubit - bit quantistico che si differenzia da un bit classico perché non è limitato a essere esclusivamente uno 0 o un 1, ma può trovarsi in una combinazione di entrambi gli stati contemporaneamente e più qubit possono essere entangled per creare una correlazione quantistica - aprendo la strada alla simulazione efficiente sia delle proprietà statiche, lo stato fondamentale, sia di quelle dinamiche, gli stati eccitati.

\slidenum{3}
Per verificare in pratica queste proprietà statiche e dinamiche, ho fissato tre obiettivi concreti. Il primo, il caso prototipo, è il dimero: due soli spin 1/2 accoppiati su 2 qubit, con l'Hamiltoniana che vedete, tre parametri fisici --- J per l'accoppiamento, b per il campo, D per l'interazione di Dzyaloshinskii--Moriya --- sufficienti a produrre una fisica non banale. Il secondo obiettivo è estendere la simulazione a tre spin interagenti, disposti a triangolo isoscele: il legame in più, J', introduce un fenomeno del tutto assente nel dimero, la frustrazione geometrica. Il terzo obiettivo è aggiungere il rumore che, in hardware reale è dovuto dall'interazione con l'ambiente.

\slidenum{4}
Il problema della dimensione dello spazio di Hilbert al crescere di N ha una risposta: l'idea risale a Feynman, nel 1982, secondo cui un computer quantistico dovrebbe poter simulare qualsiasi sistema quantistico locale - dove ogni componente del sistema interagisce solo con un numero limitato di altre componenti vicine - in modo efficiente. Lloyd, nel 1996, ha concretizzato questa intuizione usando la decomposizione di Suzuki-Trotter come strumento della dimostrazione --- la stessa che ho usato per l'evoluzione temporale. La decomposizione di Suzuki-Trotter approssima l'evoluzione temporale generata da un'Hamiltoniana somma di più termini che non commutano tra loro, dividendo il tempo totale in N intervalli piccoli e alternando gli esponenziali dei singoli termini: per N che tende a infinito la formula è esatta, per N finito introduce un errore controllabile.

\slidenum{5}
Il Variational Quantum Eigensolver (VQE) è un algoritmo ibrido quantistico-classico: la QPU (Quantum Processing Unit) prepara lo stato variazionale $|\psi(\theta)\rangle$, l'ansatz, e misura il valore di aspettazione dei vari termini dell'Hamiltoniana; la CPU (classica) li somma e decide come aggiornare i parametri $\theta$ per il giro successivo, fino a convergere verso $E_0$, l'energia dello stato fondamentale. Il Principio Variazionale assicura che il valore di aspettazione dell'energia $E(\theta)$ sarà sempre maggiore o uguale a $E_0$ da sopra. È proprio la preparazione dell'ansatz il punto su cui si gioca la qualità del risultato.

\slidenum{6}
Con lo schema generale del VQE appena visto, passiamo al caso concreto del dimero: qui confronto tre circuiti diversi, che a partire da $|00\rangle$ preparano l'ansatz, valutati sulla fidelity $F$ rispetto allo stato fondamentale vero. Il circuito hardware-efficient HA a sei parametri è il termine di confronto: una baseline generica che non usa alcuna informazione sulla fisica del problema; il circuito Physically Motivated PMA base resta dentro al sottospazio fisicamente corretto (in base al ramo di $b$); il circuito PMA esteso ha tre parametri di cui le due rotazioni più esterne rompono la simmetria.

\slidenum{7}
L'idea di PMA esteso appena vista sul dimero si estende naturalmente al trimero: qui confronto due circuiti, che a partire da $|000\rangle$ preparano l'ansatz, sempre valutati sulla fidelity $F$ rispetto allo stato fondamentale vero. RBS-2q.6 è esattamente la stessa idea del dimero portata a tre spin. Il circuito scelto W-2q.6, ha la stessa identica struttura ma sostituisce il blocco RBS con un blocco $W$ - visto nella slide precedente come blocco primitivo per il circuito PMA base del dimero - sui tre legami.

\slidenum{8}
Il termine DM, quando acceso, mescola stati con magnetizzazione $M$ diversa, rompendo la simmetria e aprendo il gap in corrispondenza del punto di incrocio - il valore di campo in cui, a $D=0$, l'energia dello stato fondamentale passa bruscamente da un livello all'altro. Sono mostrati i risultati del confronto tra gli ansatz e la loro sensibilità a questo termine. Al punto di incrocio con $b/J = 2$ e $D/J = 0.2$, HA e PMA esteso restano entrambi a fidelity 1, mentre PMA base crolla a fidelity 0.701: è il punto in cui la sovrapposizione tra settori di magnetizzazione diversa è più marcata e il circuito PMA base, confinato in un solo settore, non riesce a rappresentarla. Il punto di lavoro a destra dimostra che la degradazione dipende da $D$ anche lontano dall'incrocio: c'è sempre una mescolanza residua anche se minima.

\slidenum{9}
A sinistra, lo spettro del dimero con il termine DM acceso, i punti rossi sono l'approssimazione VQE con PMA esteso e l'accordo è praticamente perfetto su tutto il range di $b/J$. A destra, lo spettro del trimero e l'approssimazione VQE - dove è stato fatto uno zoom sul gap.

\slidenum{10}
Si inizializza il sistema fuori equilibrio (non un autostato della Hamiltoniana che implicherebbe solo un cambio di fase globale) - $|00\rangle$ per il dimero, $|000\rangle$ per il trimero - e si osserva come evolve nel tempo. I due grafici a sinistra, per il dimero e il trimero, mostrano l'evoluzione esatta a confronto con Trotter: per $t > 0$ l'Hamiltoniana non commuta con $S_z^{tot}$ a causa del termine DM, quindi lo stato evolve in una sovrapposizione di stati. Ogni stato evolve con fase diversa la loro differenza genera una interferenza dinamica che fa oscillare il segnale. L'approssimazione della evoluzione di Trotter rispetto a quella esatta scala con pendenza $-1$, mentre l'antifedeltà sullo stato specifico scala con pendenza $-2$ perché è essenzialmente il quadrato dell'errore sull'ampiezza. Sono gli scaling teorici della decomposizione di Trotter, verificati numericamente.

\slidenum{11}
Costruiamo ora un circuito quantistico con VQE e Trotter, che affianca al registro fisico un'ancilla in sovrapposizione. L'algoritmo, un Hadamard test, sfrutta l'interferenza per estrarre informazioni altrimenti inaccessibili: si prepara lo stato con VQE, si apre l'ancilla --- un qubit ausiliario senza alcun ruolo nella fisica del problema --- con una porta di Hadamard, si condiziona a essa un operatore di Pauli, si lascia evolvere il registro con Trotter, si applica un secondo Pauli anti-controllato, e infine si misura l'ancilla dopo una rotazione per la parte immaginaria o una Hadamard per la parte reale. Ogni shot dà una misura binaria; sommandone molti si ottengono le probabilità sperimentali che ricostruiscono parte reale e immaginaria del correlatore. Il grafico a destra confronta la curva esatta, Trotter senza rumore e Trotter con rumore a $N=20$, calibrato sui parametri mediani di ibm\_torino. - VQE con NoiseModel + Trotter rumoroso + readout rumoroso.

\slidenum{12}
I due grafici a sinistra, per dimero e trimero, mostrano che senza rumore di readout - solo rumore di gate reale in Trotter e VQE ottimizzato senza rumore inizialmente e rieseguito agganciando il NoiseModel costruito con i parametri di calibrazione - la fedeltà cresce fino a un massimo, $N^* = 8$ per il dimero e $N^* = 3$ per il trimero, per poi decrescere: è il compromesso tra errore di Trotter, che scende con $N$, e rumore accumulato, che sale con $N$. I grafici centrali mostrano la stessa cosa ma con il VQE con parametri ottimizzati su circuito rumoroso. Questo mostra che riottimizzare vedendo il rumore non porta alcun vantaggio misurabile: le curve restano praticamente sovrapposte e $N^*$ resta lo stesso in entrambi i casi. A destra, il confronto diretto dimero contro trimero mostra che con la stessa calibrazione di rumore, la maggiore profondità del circuito nel trimero produce una degradazione di fedeltà quasi doppia.

\slidenum{13}
Ho studiato sistemi modello di spin interagenti come banco di prova per algoritmi di simulazione quantistica. Lo stato fondamentale è determinato con il VQE, un algoritmo ibrido che fa lavorare insieme hardware classico e quantistico. Le performance cambiano quando si introduce il rumore realistico dell'hardware, e l'ottimizzazione ne risente. Inoltre sfruttare le simmetrie del sistema, invece di un ansatz generico, riduce i gate e quindi l'errore e quindi migliora le performance.

\end{document}
